% GENERATED FILE. Edit canonical lessons, then run npm run generate:books.
% canonical-source-hash: fnv1a64:2fa585db33b5d669
% canonical-lessons: KA-C472-halagu, KA-C472-emaru, KA-C472-paragu, KA-C472-asraya, KA-C472-sahayavani

\chapter{When Things Go Wrong}
\label{ch:ka-when-things-go-wrong}
\hlchaptermodality{472}

\begin{chapteropening}
\textbf{By the end of this chapter:} \emph{I can say to get damaged, to be off guard, to escape unharmed, shelter and a helpline.}

Complete the last lesson of chapter 472: I can say to get damaged, to be off guard, to escape unharmed, shelter and a helpline.
\end{chapteropening}

\section[\texorpdfstring{\kn{ಹಾಳಾಗು} (hāḷāgu)}{ಹಾಳಾಗು (hāḷāgu)}]{\kn{ಹಾಳಾಗು} (hāḷāgu) — to get damaged; to be ruined}

\label{lesson:KA-C472-halagu}



\begin{quote}
Before the new one: say the Kannada for to tug, then the Kannada for to spill over.
\end{quote}

\subsection*{You'll want to know: \kn{ಹಾಳಾಗು}}
\textbf{\kn{ಹಾಳಾಗು}} — \emph{hāḷāgu} — \textquotedblleft{}to get damaged; to be ruined\textquotedblright{}.

\begin{quote}
\kn{ಮಳೆಯಲ್ಲಿ} \kn{ನನ್ನ} \kn{ಫೋನ್} \kn{ಹಾಳಾಯಿತು}.
\end{quote}

\begin{quote}
\emph{maḷeyalli nanna phōn hāḷāyitu.}
\end{quote}

\begin{quote}
My phone got ruined in the rain.
\end{quote}

Past (it …): \kn{ಹಾಳಾಯಿತು}, \emph{hāḷāyitu}.

\kn{ಹಾಳು} \kn{ಮಾಡು}, to spoil something.

\begin{cousinweb}
\textbf{\kn{ಹಾಳಾಗು}} (\emph{hāḷāgu}) is native: \kn{ಹಾಳು}, \textquotedblleft{}ruin, a waste\textquotedblright{} (DEDR 4110, Kannada \emph{pāḻ}), joined with \kn{ಆಗು}, to become.
\end{cousinweb}

\begin{grammarlens}[title={What you've built}]
One more word for getting around.
\end{grammarlens}

\subsection*{Your turn}
\begin{itemize}
\raggedright
  \item \emph{Say it:} \emph{hāḷāgu}
  \item \emph{Say it:} \emph{hāḷāgu}, once more
  \item \emph{Say it:} \emph{maḷeyalli nanna phōn hāḷāyitu.} — My phone got ruined in the rain.
  \item \emph{Recall:} say the Kannada for to tug, then the Kannada for to spill over, then say \emph{hāḷāgu} again
\end{itemize}

\begin{tcolorbox}[breakable,colback=teal!4,colframe=teal!35!black,title={Before you move on}]
What does \kn{ಹಾಳಾಗು} mean? (\textquotedblleft{}To get damaged; to be ruined\textquotedblright{}.) Say it once more.
\end{tcolorbox}
\section[\texorpdfstring{\kn{ಏಮಾರು} (ēmāru)}{ಏಮಾರು (ēmāru)}]{\kn{ಏಮಾರು} (ēmāru) — to be off guard; to be fooled}

\label{lesson:KA-C472-emaru}



\begin{quote}
Before the new one: say the Kannada for to spill over, then the Kannada for to get damaged.
\end{quote}

\subsection*{You'll want to know: \kn{ಏಮಾರು}}
\textbf{\kn{ಏಮಾರು}} — \emph{ēmāru} — \textquotedblleft{}to be off guard; to be fooled\textquotedblright{}.

\begin{quote}
\kn{ನಾನು} \kn{ಒಂದು} \kn{ಕ್ಷಣ} \kn{ಏಮಾರಿದೆ}, \kn{ಆಗ} \kn{ಯಾರೋ} \kn{ನನ್ನ} \kn{ಚೀಲ} \kn{ಕದ್ದರು}.
\end{quote}

\begin{quote}
\emph{nānu ondu kṣaṇa ēmāride, āga yārō nanna cīla kaddaru.}
\end{quote}

\begin{quote}
I let my guard down for a moment, and then someone stole my bag.
\end{quote}

Past (I …): \kn{ಏಮಾರಿದೆ}, \emph{ēmāride}.

A moment of carelessness that someone takes advantage of.

\begin{cousinweb}
\textbf{\kn{ಏಮಾರು}} (\emph{ēmāru}) is usually linked to the native Dravidian family that DEDR 898 collects, Tamil and Telugu \emph{ēmāṟu}, \textquotedblleft{}to be beguiled, to be deceived\textquotedblright{}; DEDR gives the Kannada only with a query, as \emph{vēmāṟu}, \textquotedblleft{}to be careless\textquotedblright{}, so the link is not settled.
\end{cousinweb}

\begin{grammarlens}[title={What you've built}]
One more word for getting around.
\end{grammarlens}

\subsection*{Your turn}
\begin{itemize}
\raggedright
  \item \emph{Say it:} \emph{ēmāru}
  \item \emph{Say it:} \emph{ēmāru}, once more
  \item \emph{Say it:} \emph{nānu ondu kṣaṇa ēmāride, āga yārō nanna cīla kaddaru.} — I let my guard down for a moment, and then someone stole my bag.
  \item \emph{Recall:} say the Kannada for to spill over, then the Kannada for to get damaged, then say \emph{ēmāru} again
\end{itemize}

\begin{tcolorbox}[breakable,colback=teal!4,colframe=teal!35!black,title={Before you move on}]
What does \kn{ಏಮಾರು} mean? (\textquotedblleft{}To be off guard; to be fooled\textquotedblright{}.) Say it once more.
\end{tcolorbox}
\section[\texorpdfstring{\kn{ಪಾರಾಗು} (pārāgu)}{ಪಾರಾಗು (pārāgu)}]{\kn{ಪಾರಾಗು} (pārāgu) — to escape unharmed; to get through safely}

\label{lesson:KA-C472-paragu}



\begin{quote}
Before the new one: say the Kannada for to get damaged, then the Kannada for to be off guard.
\end{quote}

\subsection*{You'll want to know: \kn{ಪಾರಾಗು}}
\textbf{\kn{ಪಾರಾಗು}} — \emph{pārāgu} — \textquotedblleft{}to escape unharmed; to get through safely\textquotedblright{}.

\begin{quote}
\kn{ಅಪಘಾತದಲ್ಲಿ} \kn{ಎಲ್ಲಾ} \kn{ಪ್ರಯಾಣಿಕರು} \kn{ಪಾರಾದರು}.
\end{quote}

\begin{quote}
\emph{apaghātadalli ellā prayāṇikaru pārādaru.}
\end{quote}

\begin{quote}
All the passengers came through the accident safely.
\end{quote}

Past (I …): \kn{ಪಾರಾದೆ}, \emph{pārāde}.

What one escapes takes -\kn{ಇಂದ}: \kn{ಅಪಾಯದಿಂದ} \kn{ಪಾರಾಗು}.

\begin{cousinweb}
\textbf{\kn{ಪಾರಾಗು}} (\emph{pārāgu}) joins Sanskrit \emph{pāra}, \textquotedblleft{}the further bank, the far end\textquotedblright{} (CDIAL 8100), and the native verb \kn{ಆಗು}: to reach the other side.
\end{cousinweb}

\begin{grammarlens}[title={What you've built}]
One more word for getting around.
\end{grammarlens}

\subsection*{Your turn}
\begin{itemize}
\raggedright
  \item \emph{Say it:} \emph{pārāgu}
  \item \emph{Say it:} \emph{pārāgu}, once more
  \item \emph{Say it:} \emph{apaghātadalli ellā prayāṇikaru pārādaru.} — All the passengers came through the accident safely.
  \item \emph{Recall:} say the Kannada for to get damaged, then the Kannada for to be off guard, then say \emph{pārāgu} again
\end{itemize}

\begin{tcolorbox}[breakable,colback=teal!4,colframe=teal!35!black,title={Before you move on}]
What does \kn{ಪಾರಾಗು} mean? (\textquotedblleft{}To escape unharmed; to get through safely\textquotedblright{}.) Say it once more.
\end{tcolorbox}
\section[\texorpdfstring{\kn{ಆಶ್ರಯ} (āśraya)}{ಆಶ್ರಯ (āśraya)}]{\kn{ಆಶ್ರಯ} (āśraya) — shelter; refuge}

\label{lesson:KA-C472-asraya}



\begin{quote}
Before the new one: say the Kannada for to be off guard, then the Kannada for to escape unharmed.
\end{quote}

\subsection*{You'll want to know: \kn{ಆಶ್ರಯ}}
\textbf{\kn{ಆಶ್ರಯ}} — \emph{āśraya} — \textquotedblleft{}shelter; refuge\textquotedblright{}.

\begin{quote}
\kn{ಮಳೆ} \kn{ಬಂದಾಗ} \kn{ನಾವು} \kn{ಒಂದು} \kn{ಅಂಗಡಿಯಲ್ಲಿ} \kn{ಆಶ್ರಯ} \kn{ಪಡೆದೆವು}.
\end{quote}

\begin{quote}
\emph{maḷe bandāga nāvu ondu aṅgaḍiyalli āśraya paḍedevu.}
\end{quote}

\begin{quote}
When it rained we took shelter in a shop.
\end{quote}

\kn{ಆಶ್ರಯ} \kn{ಪಡೆ}, to take shelter; \kn{ಆಶ್ರಯ} \kn{ಕೊಡು}, to give shelter.

\begin{cousinweb}
\textbf{\kn{ಆಶ್ರಯ}} (\emph{āśraya}) is borrowed from Sanskrit \emph{āśraya}, \textquotedblleft{}a resting place, a refuge, a dwelling\textquotedblright{} (CDIAL 1466).
\end{cousinweb}

\begin{grammarlens}[title={What you've built}]
One more word for getting around.
\end{grammarlens}

\subsection*{Your turn}
\begin{itemize}
\raggedright
  \item \emph{Say it:} \emph{āśraya}
  \item \emph{Say it:} \emph{āśraya}, once more
  \item \emph{Say it:} \emph{maḷe bandāga nāvu ondu aṅgaḍiyalli āśraya paḍedevu.} — When it rained we took shelter in a shop.
  \item \emph{Recall:} say the Kannada for to be off guard, then the Kannada for to escape unharmed, then say \emph{āśraya} again
\end{itemize}

\begin{tcolorbox}[breakable,colback=teal!4,colframe=teal!35!black,title={Before you move on}]
What does \kn{ಆಶ್ರಯ} mean? (\textquotedblleft{}Shelter; refuge\textquotedblright{}.) Say it once more.
\end{tcolorbox}
\section[\texorpdfstring{\kn{ಸಹಾಯವಾಣಿ} (sahāyavāṇi)}{ಸಹಾಯವಾಣಿ (sahāyavāṇi)}]{\kn{ಸಹಾಯವಾಣಿ} (sahāyavāṇi) — a helpline}

\label{lesson:KA-C472-sahayavani}



\begin{quote}
Before the new one: say the Kannada for to escape unharmed, then the Kannada for shelter.
\end{quote}

\subsection*{You'll want to know: \kn{ಸಹಾಯವಾಣಿ}}
\textbf{\kn{ಸಹಾಯವಾಣಿ}} — \emph{sahāyavāṇi} — \textquotedblleft{}a helpline\textquotedblright{}.

\begin{quote}
\kn{ಚೀಲ} \kn{ಸಿಗದಿದ್ದಾಗ} \kn{ನಾವು} \kn{ಸಹಾಯವಾಣಿಗೆ} \kn{ಫೋನ್} \kn{ಮಾಡಿದೆವು}.
\end{quote}

\begin{quote}
\emph{cīla sigadiddāga nāvu sahāyavāṇige phōn māḍidevu.}
\end{quote}

\begin{quote}
When we could not find the bag we rang the helpline.
\end{quote}

\kn{ಸಹಾಯ}, help, plus \kn{ವಾಣಿ}, voice.

\begin{cousinweb}
\textbf{\kn{ಸಹಾಯವಾಣಿ}} (\emph{sahāyavāṇi}) is a modern compound built from Sanskrit words: \emph{sahāya}, \textquotedblleft{}a companion, a helper\textquotedblright{} (CDIAL 13309), and \emph{vāṇī}, \textquotedblleft{}voice\textquotedblright{} (CDIAL 11490).
\end{cousinweb}

\begin{grammarlens}[title={What you've built}]
One more word for getting around.
\end{grammarlens}

\subsection*{Your turn}
\begin{itemize}
\raggedright
  \item \emph{Say it:} \emph{sahāyavāṇi}
  \item \emph{Say it:} \emph{sahāyavāṇi}, once more
  \item \emph{Say it:} \emph{cīla sigadiddāga nāvu sahāyavāṇige phōn māḍidevu.} — When we could not find the bag we rang the helpline.
  \item \emph{Recall:} say the Kannada for to escape unharmed, then the Kannada for shelter, then say \emph{sahāyavāṇi} again
\end{itemize}

\begin{tcolorbox}[breakable,colback=teal!4,colframe=teal!35!black,title={Before you move on}]
What does \kn{ಸಹಾಯವಾಣಿ} mean? (\textquotedblleft{}A helpline\textquotedblright{}.) Say it once more.
\end{tcolorbox}
